
\chapter{מבוא לחילוץ אותות סינוסואידליים}

\section{אתגר חילוץ אותות מרעש}

אחת הבעיות הקלאסיות בעיבוד אותות היא חילוץ מרכיבים סינוסואידליים טהורים מאותות מעורבבים המכילים רעש (\textenglish{noisy mixed signals}). בעולם האמיתי, מדידות פיזיות כמו אותות ביו-רפואיים, תדרי רדיו ומוסיקה דיגיטלית מורכבות מעירוב של מספר מרכיבי תדר, כאשר כל אחד מהם מוסיף רעש מתחום גאוסי (\textenglish{Gaussian noise}) או אימפולסיבי (\textenglish{impulsive noise}). הפרדת הרכיבים השונים דורשת גישה מתוחכמת המסוגלת ללמוד מבנים עמוקים בנתונים.

גישות קלאסיות לבעיה זו כוללות פירוק \textenglish{Fourier} (\textenglish{FFT}) ופילטרים דיגיטליים כמו \textenglish{Wiener filter}. אולם שיטות אלו מניחות מידה גבוהה של ידע מוקדם על התדרים הנכללים ואינן מסוגלות להתאים את עצמן לשינויים דינמיים בתדר, משרעת או פאזה. רשתות נוירונים עמוקות, ובמיוחד ארכיטקטורות רקורנטיות (\textenglish{RNN}) ורשתות \textenglish{LSTM} (\textenglish{Long Short-Term Memory}) \cite{hochreiter1997lstm}, הציגו יכולות מרשימות בלמידה אוטומטית של מבנים זמניים בנתונים.

\section{יישומים מעשיים}

לחילוץ גלי סינוס מרעש יש יישומים מגוונים: בתחום הרפואה, ניתוח אותות \textenglish{ECG} ו-\textenglish{EEG} דורש הפרדה של מרכיבים פיזיולוגיים מרעש חיישנים \cite{valentini2019speech}. בתחום התקשורת, זיהוי תדרים בסביבות עם הפרעה דורש עמידות גבוהה לרעש. מחקר זה מציג מסגרת למידה עמוקה לחילוץ פרמטרי גלי סינוס: משרעת, תדר ופאזה, עם ביצועים עדיפים על גישות מסורתיות.

\section{תרומות המאמר}

מאמר זה מציג שלוש תרומות עיקריות: ראשית, ארכיטקטורת \textenglish{BiLSTM Encoder-Decoder} ייעודית לחילוץ פרמטרי; שנית, פונקציית אובדן מבוססת \textenglish{SI-SNR} (\textenglish{Scale-Invariant Signal-to-Noise Ratio}) שהוכחה יציבה יותר מ-\textenglish{MSE} בתנאי רעש חזק; ושלישית, הדגמת שיפור של עד \textenglish{4 dB} ב-\textenglish{SI-SNR} על פני קו הבסיס הנוכחי \cite{luo2019convtasnet} בניסויי \textenglish{benchmark} מקובלים.
